\chapter{Physical Interpretation and Unresolved Comparisons}
\label{ch:assessment}

\section{Mathematical construction and physical validation}

A coherent mathematical construction, a controlled physical model, and an empirically supported theory are different kinds of result. The finite-depth covariant and canonical identities concern a stated model. The production spectrum additionally requires a controlled solution of the coupled dynamics. Empirical identification requires comparison with data and competing models.

The energy-budget obstruction locates a failure in the specified production history. It does not invalidate projector reconstruction, matrix-corner calculus, or the conditional transport of a prescribed distribution. It does prevent that distribution from being called a prediction of the microscopic crossing without a new dynamical argument.

\section{Established model statements}

The central results separate symmetry, conservation, approximation, and physical production.

\begin{longtable}{@{}P{0.36\textwidth}P{0.20\textwidth}P{0.34\textwidth}@{}}
\toprule
Result & Status & Scientific role\\
\midrule
Level-orthogonal quadratic operator classification & theorem & isolates one coefficient per refinement level\\
Branch-homogeneity recurrence & theorem & reduces those coefficients to one geometric family\\
Leading analytic spectral interaction & theorem within stated parity, locality, and exact-order assumptions & explains one channel but leaves an allowed order-zero coupling\\
Covariant stress-exchange identity & theorem & supplies an exact conservation diagnostic\\
Total hypersurface algebra with scalar fiber & theorem & shows that the finite-depth matter sector is canonically compatible after gravity is included\\
Resolved matter $H$--$H$ test and negative controls & controlled diagnostic & separates spectral headroom from aliasing and omitted-gradient failures\\
Arbitrary-order continuum moments & theorem for the interpolated density model & extends the production ratios beyond \(\sigma=1\)\\
Fixed-prime Mellin decomposition & theorem & corrects the interpretation of the continuum approximation\\
Specified benchmark energy audit & controlled diagnostic & proves the prescribed crossing is not weakly backreacting\\
Finite-depth one-loop power counting & theorem at leading ultraviolet order & exposes the need for a cutoff or cancellation\\
Heavy-modulus realization & model proposition & supplies a testable route to a common scale multiplier\\
\bottomrule
\end{longtable}

The energy-budget result is a central obstruction.  The specified rapid-crossing spectrum is not a small perturbation of its background.  Its unprojected energy exceeds the declared reservoir by \(\BridgeOverdraw\).  A free crossing followed by a common normalization factor therefore cannot be read as a self-consistent production history.

This negative result is constructive.  It localizes the failure upstream of the late transport calculation and gives exact equations for replacing the projection.  It also prevents further phenomenology from accumulating on an uncontrolled microscopic input.

\section{Limits of the model}

The following claims remain outside the proved domain.
\begin{enumerate}[leftmargin=2.2em]
\item The projector axioms do not select \(p\), \(\sigma\), \(\Lambda\), \(m_b\), \(\gamma\), the potential parameters, or the initial state.
\item The symmetry classification does not calculate the numerical value \(\gamma=16\).
\item The order-one backreaction threshold is not a solution of the coupled background and mode equations or an energy-partition equation.
\item The projected Gaussian ratios are not preserved under strong backreaction unless a dynamical calculation demonstrates that fact.
\item The ordinary infinite bosonic tower is not radiatively finite under the finite-depth continuation studied here.
\item The modulus completion does not solve the cosmological constant problem or prove ultraviolet completeness.
\item The internal verification archive is not independent reproduction, peer review, or experimental evidence.
\item No dark matter or dark energy identification follows solely from a match to a prescribed background abundance.
\end{enumerate}

\section{A controlled production comparison}

An additional observable template does not resolve the upstream production problem.  It is a finite-depth backreacting production calculation with one conserved renormalized stress tensor.  The system in \eqref{eq:back-hartree-phi}--\eqref{eq:back-hartree-friedmann} is the minimum viable model.

The calculation should be organized as a sequence of gates.

\begin{longtable}{@{}P{0.13\textwidth}P{0.35\textwidth}P{0.42\textwidth}@{}}
\toprule
Gate & Required result & Failure meaning\\
\midrule
G1 & Renormalized total continuity residual converges to zero with time step and momentum resolution & implementation or subtraction is inconsistent\\
G2 & Observables stabilize across at least three finite depths and two momentum cutoffs & prediction is regulator dominated\\
G3 & The produced spectrum is obtained from the coupled evolution without a common post hoc projection & microscopic production is actually solved\\
G4 & Interaction and loop expansion parameters are below declared bounds, or a nonperturbative method is used & free-mode interpretation is invalid\\
G5 & A viable region reproduces abundance, pressure, and coldness simultaneously & the proposed cosmological branch survives\\
G6 & At least one linked observable differs from a simpler model after nuisance parameters are fit & the construction becomes falsifiable rather than only flexible\\
\bottomrule
\end{longtable}

\begin{negativebox}[Go or revise]
If no cutoff-stable solution passes G1--G5, the specified cosmological branch should be revised or retired.  The projector, corner, and constrained-system results would remain as separate mathematical work.  A failed dark-sector realization would not invalidate them.
\end{negativebox}

\section{Minimum computational experiment}

A credible first run does not need the full infinite tower.  It needs a converged finite-depth sequence.

\begin{enumerate}[leftmargin=2.2em]
\item Choose \(H=1,2,3,4\) and publish the physical largest mass at every depth.
\item Initialize fourth-order adiabatic vacuum modes over a documented comoving momentum grid.
\item Use the same adiabatic subtraction terms in \(Y_H^{\mathrm{ren}}\), \(\rho_H^{\mathrm{ren}}\), and \(p_H^{\mathrm{ren}}\).
\item Evolve \(\phi\), all retained modes, and the Friedmann equation together through the crossing.
\item Monitor the covariant identity in \eqref{eq:cov-ward} and the FLRW exchanges in \eqref{eq:cov-exchange-chi}--\eqref{eq:cov-exchange-phi} at every step.
\item Repeat with doubled time resolution, doubled momentum resolution, and shifted ultraviolet cutoffs.
\item Report the unprojected spectrum, energy, pressure, rest fraction, level occupancy, and their convergence envelopes.
\item Scan \(\lambda\), \(\gamma\), \(\Lambda/\mu\), \(m_b/\mu\), and the invariant crossing condition.  Do not use a coordinate value of the scale factor as the only clock.
\end{enumerate}

The first scientific deliverable from this experiment should be a survival map, not a best-fit point.  Its axes should distinguish the weakly coupled, strongly backreacting, tachyonic, cutoff-sensitive, and radiatively unstable regions.

\section{A second foundational track}

The strongest purely gravitational continuation is different.  It is an end-to-end nonlinear Gowdy corner--BFV continuum bridge.  That track should derive the ungauge-fixed \(T^2\)-symmetric sector from the real polynomial Lorentzian parent, construct an explicit corner derivative and variable-shift generator, instantiate a cyclic contraction, prove finite BFV nilpotency and refinement-path equivalence, and establish uniform local-in-time convergence after areal gauge on \([t_0,T]\) with \(t_0>0\).

That theorem would close the gap between the abstract matrix-corner calculus and the filtered Fourier realizations used in the later gravitational chapters.  It is the right foundational milestone alongside the matter comparison.  It should not be replaced by an unrestricted \(3+1\) or quantum claim before the symmetry-reduced continuum map is explicit.

\section{Decision table for the program}

\begin{longtable}{@{}P{0.27\textwidth}P{0.31\textwidth}P{0.32\textwidth}@{}}
\toprule
Outcome & Interpretation & Next action\\
\midrule
Backreacting spectrum remains close to the prescribed ansatz in a controlled region & cosmological bridge gains microscopic support & rerun late observables with propagated theory uncertainty\\
Spectrum survives but is substantially reshaped & mechanism may survive, benchmark does not & replace frozen initial distribution and refit every downstream stage\\
Only strongly coupled solutions exist & free production calculation is the wrong tool & formulate and validate a nonperturbative treatment\\
Finite-depth results fail to stabilize & tower prediction is regulator dependent & identify a physical endpoint or ultraviolet cancellation\\
No viable production region exists & dark-sector realization is falsified in this form & retain mathematical framework and revise the matter completion\\
Gowdy bridge succeeds independently & foundational architecture gains a nonlinear continuum theorem & extend cautiously to additional symmetry classes\\
\bottomrule
\end{longtable}

\section{Final claim boundary}

The construction supplies exact internal identities, conditional gravitational statements, and an explicitly testable fiber model. Its specified cosmological production benchmark fails a quantified energy-budget test. A viable physical realization must control backreaction, ultraviolet sensitivity, and the transport of any resulting uncertainty through the observable calculations.

The cosmology chapter therefore begins from a prescribed normalized distribution. Its deductions retain their conditional meaning. They do not assume that the unresolved production comparison has succeeded.
